% 三个主张层次及同层职责对照；横线没有因果或逻辑推出含义。
\begin{tikzpicture}[figure picture,foundation picture,
 font=\figurefont\fontsize{8.25}{10.5}\selectfont,
 trust cell/.style={draw=FigureStroke,line width=.7pt,rounded corners=1.5pt},
 trust note/.style={font=\FigureNoteFont\fontsize{7.75}{10}\selectfont,align=left}]
 \path[use as bounding box] (0,-2) rectangle (169,65);
 \node at (20,59) {主张层次};
 \node at (77,59) {关键问题};
 \node at (141.5,59) {责任分工};
 \foreach \y/\fillcol/\stage in {44/FigureBlueFill/属性与规范,26/MutedGreenFill/证据能力,8/MutedPurpleFill/失效与处置}{
  \draw[trust cell,fill=\fillcol] (0,{\y-7}) rectangle (40,{\y+7});
  \draw[trust cell,fill=\fillcol] (49,{\y-7}) rectangle (105,{\y+7});
  \draw[trust cell,fill=\fillcol] (114,{\y-7}) rectangle (169,{\y+7});
  \node at (20,\y) {\stage};
  \draw[draw=FigureStroke,line width=.65pt] (40,\y)--(49,\y) (105,\y)--(114,\y);
 }
 \node[trust note] at (77,44) {必须满足什么？\\适用于谁、什么范围？};
 \node[trust note] at (77,26) {结论有多强？\\还缺什么证据？};
 \node[trust note] at (77,8) {何时失效？\\何时限制、回退？};
 \node[trust note] at (141.5,44) {任务责任方设定要求\\结合受影响者与领域规则};
 \node[trust note] at (141.5,26) {技术人员论证\\复核者检查边界};
 \node[trust note] at (141.5,8) {有权主体决策\\运行责任方处置};
\end{tikzpicture}
